% esl-crossbar-array — Parametric nrows x ncols VMM crossbar-cell array (contract §1a crossbar-array), generated by nested \foreach loops with row/col-indexed node names, row voltage sources, and per-column KCL summation nodes.
% THESIS: Row and column wiring is drawn as ONE colinear rail per row/column, never point-to-point per cell, so the array stays both correct-looking and legible at scale -- proven at 64x64.
% Concept grafted from: esl-crossbar-kcl (cp-4856 PoC) crossbar/KCL role and color assignments.
% Intent record: intent.yaml (contract §4 shape; ADR-0005 D3 §4 row).
\documentclass[border=8pt]{standalone}

% --- packages (mirror these in template.meta.json "requires") ---
\usepackage[dvipsnames]{xcolor}
\usepackage{tikz}
\usetikzlibrary{calc, arrows.meta}

% --- palette: ESL domain-semantic tokens (contract §2, ADR-0005 D7 -- extends
% opentikz's ot* palette, does not remap onto it; see
% reference/color-palettes/color-palettes.md "ESL domain-semantic tokens").
% Same token bindings as esl-crossbar-kcl (inputdomain=blue, weightdomain=green,
% analogdomain=orange, digitaldomain=gray) so a \foreach-generated array reads
% as the same family, not a different figure.
\colorlet{inputdomain}{blue}     % inputs / activations (row voltage sources)
\colorlet{weightdomain}{green}   % analog weight cells
\colorlet{analogdomain}{orange}  % analog current / accumulation (KCL, rails)
\colorlet{digitaldomain}{gray}   % neutral axis/index annotation

\begin{document}
% ==== parameters (edit these; this is the \foreach generator's parameter surface —
% nothing below this block should ever be edited to change array shape/size) =====
\def\nrows{2}            % number of crossbar rows (VMM input rows)
\def\ncols{2}             % number of crossbar columns (VMM output columns)
\def\rowgap{0.7}          % vertical gap between rows (cm)
\def\colgap{0.7}          % horizontal gap between columns (cm)
\def\cellsize{0.5}        % crossbar-cell node size (cm)
\def\celllabelmode{full}  % full|none -- legibility control for per-cell $G_{r,c}$ text
\def\showkcl{1}           % 1 = draw per-column KCL summation node; 0 = omit
\def\vlabelprefix{V}      % row voltage-source label prefix
\def\edgelabelstride{1}   % show a row/col index label only every Nth line (1 = every line)
% =========================================================================

\begin{tikzpicture}[
    font=\small,
    >=Stealth,
    cell/.style={draw, rectangle, thick, fill=weightdomain!15,
                 minimum size=\cellsize cm, inner sep=0pt, font=\tiny},
    vsource/.style={draw=inputdomain!55, circle, line width=0.8pt, fill=inputdomain!18,
                    minimum size=\cellsize cm, inner sep=0pt,
                    font=\tiny\itshape},
    kclnode/.style={draw, circle, thick, fill=analogdomain!35,
                    minimum size=\cellsize cm, inner sep=0pt, font=\tiny},
    awire/.style={analogdomain!80!black, thick},
    vwire/.style={inputdomain!60!black, thick},
    axislbl/.style={font=\tiny, digitaldomain!40!black},
  ]

  % legibility sentinel: \celllabelmode is compared against this literal.
  \def\CellShowFull{full}

  % ==== generator: nrows x ncols crossbar-cell array, row-indexed voltage
  % sources, column-indexed KCL nodes (contract §1a crossbar-array: "a full
  % nrows x ncols VMM cell array, placed by row/col index") ====
  \foreach \r in {1,...,\nrows} {
    \pgfmathsetmacro{\py}{-(\r-1)*\rowgap}
    \node[vsource] (vin-\r) at (0,\py) {};
    \pgfmathtruncatemacro{\atstride}{mod(\r-1,\edgelabelstride)}
    \ifnum\atstride=0
      \node[axislbl, anchor=east] at ([xshift=-2pt]vin-\r.west) {\vlabelprefix\r};
    \else
      \ifnum\r=\nrows
        \node[axislbl, anchor=east] at ([xshift=-2pt]vin-\r.west) {\vlabelprefix\r};
      \fi
    \fi
    \foreach \c in {1,...,\ncols} {
      \pgfmathsetmacro{\px}{\c*\colgap}
      \node[cell] (cell-\r-\c) at (\px,\py) {%
        \ifx\celllabelmode\CellShowFull $G_{\r,\c}$\fi
      };
    }
    % row rail: one draw from the voltage source through every cell in the row
    % (all colinear at y=\py) -- O(1) edges per row, not O(cols) point-to-point.
    \draw[vwire] (vin-\r) -- (cell-\r-\ncols);
  }

  % ==== column rails + optional KCL summation nodes (contract §1a kcl-node:
  % "summation node collecting column current") ====
  \pgfmathsetmacro{\kclrow}{-(\nrows-1)*\rowgap-\rowgap}
  \foreach \c in {1,...,\ncols} {
    \pgfmathsetmacro{\px}{\c*\colgap}
    \ifnum\showkcl=1
      \node[kclnode] (kcl-\c) at (\px,\kclrow) {$\Sigma$};
      % column rail: one draw from the top cell through every cell in the
      % column down to its KCL node (all colinear at x=\px).
      \draw[awire] (cell-1-\c) -- (kcl-\c);
      \def\colbottom{kcl-\c}
    \else
      \draw[awire] (cell-1-\c) -- (cell-\nrows-\c);
      \def\colbottom{cell-\nrows-\c}
    \fi
    \pgfmathtruncatemacro{\atstride}{mod(\c-1,\edgelabelstride)}
    \ifnum\atstride=0
      \node[axislbl, anchor=north] at ([yshift=-2pt]\colbottom.south) {\c};
    \else
      \ifnum\c=\ncols
        \node[axislbl, anchor=north] at ([yshift=-2pt]\colbottom.south) {\c};
      \fi
    \fi
  }

\end{tikzpicture}
\end{document}
